\chapter*{Posfácio\smallskip\subtitulo{A interminável busca por um lugar no mundo}}
\addcontentsline{toc}{chapter}{Posfácio, \textit{por Paulina Cho}}

\begin{flushright}
\textsc{paulina cho}
\end{flushright}

\noindent{}Percorrer \emph{``Perspectivas judaico-diaspóricas sobre o lugar''} é
testemunhar um povo redefinindo e santificando o próprio conceito de
lar. É uma exploração de como, diante da dispersão, as noções abstratas
de memória e ritual se convertem em tecnologia de enraizamento para
forjar um território. A própria obra textual também é um gesto de
inscrição no tempo, uma tapeçaria, um mapa, um lugar.

Crespin revela que lugar não é geografia. É a forma como uma comunidade
decide, contra todas as probabilidades, continuar existindo. Seja por
meio da santificação de um texto, da observância de um ritual ou do
simples ato de resistir ao apagamento. O lugar torna-se verbo, prática,
repetição ritual de pertencimento. Assim, criar um "lugar" é um ato
contínuo e desafiador, não exige permanência, exige recolher o que sobra
das cinzas e insistir que ali ainda há algo que pode florescer.

A construção de lar identitário emerge das táticas diaspóricas mais
íntimas e cotidianas.

Minha mãe, que saiu da Coreia em 1985, criou seu lugar nas cozinhas dos
apartamentos em que moramos em São Paulo. Talvez a culinária seja essa
primeira tradução que fazemos ao chegar a um novo território. Adaptando
ingredientes locais para recriar seus temperos e conservas, ela foi
descobrindo, desenvolvendo e dominando uma nova linguagem sensorial. Ela
traduziu, a seu modo, os fazeres, os costumes e a forma de se comunicar
com o mundo que a cerca.

A minha mãe encarna aquilo que Gloria Anzaldúa chamaria de existir na
fronteira, condição de estar entre, que se torna, ela mesma, um lugar de
identidade. Entre-códigos próprios, ao adaptar técnicas ancestrais a
ingredientes brasileiros, entre-lugares ao transformar a cozinha
paulistana em território de memória coreana.

Como nos lembra Stephanie Borges:

\begin{quote}
"A diáspora não nos oferece um refúgio; no entanto, nada nos impede de
criar refúgios, santuários, recantos".
\end{quote}

Mas e quando as cinzas ainda estão quentes? Quando o fogo não parou de
queimar?

Hoje é 13 de junho de 2025 e tenho medo de me conectar ao mundo. De
abrir as redes sociais, ligar a TV, falar com os amigos. As notícias
chegam mesmo assim. Dou de cara com a brutal realidade de múltiplos
conflitos que se desenrolam ao nosso redor.

Se a condição diaspórica, como argumenta Crespin, consiste em
transformar o trânsito em pertencimento, hoje o trânsito parece perpétuo
para muitos, sendo esta sua condição de estar no mundo.

Penso na minha família, que cruzou oceanos carregando apenas histórias e
a esperança de que, em algum lugar --- qualquer lugar ---, pudesse
plantar raízes. Eles sabiam algo sobre criar a ideia de lugar onde esse
lugar não existe. Mas isso era antes. A violência evoluiu. Agora ela não
apenas tira pessoas de lugares, ela mata os próprios lugares definidos
por Crespin. A destruição se dá tanto na geografia quanto na própria
possibilidade de cosmogonizar novos espaços. Estamos testemunhando a
destruição da continuidade cultural, da memória coletiva, dos
fundamentos da identidade. Uma violência que é tanto física quanto
psicológica e leva à completa desumanização daquele que é condenado a
partir. O apagamento deliberado de uma paisagem que abriga o
pertencimento e a possibilidade de retorno. O objetivo da distribuição é
a fragmentação a ponto de inviabilizar a soberania futura e, assim,
redefinir permanentemente a composição.

A Guerra Russo-Ucraniana tornou-se uma luta devastadora pela identidade
nacional e integridade territorial, com milhões de deslocados e a
contestação do próprio mapa de uma nação. Na Síria, uma teia de facções
e a intervenção de várias potências estrangeiras criaram uma paisagem
fraturada. Conflitos no Sudão, na República Democrática do Congo e em
Mianmar são lembretes angustiantes de como o controle do território está
interligado à própria possibilidade de sobrevivência das comunidades. As
tensões no Mar da China Meridional evidenciam como os conflitos
geopolíticos transformam extensões marítimas em zonas de disputa. Em Los
Angeles, agentes federais invadem casas nos bairros majoritariamente
habitados por imigrantes, separando famílias sem aviso prévio. Em São
Paulo, a comunidade da Favela do Moinho enfrenta remoção forçada. Um
estilhaçamento cotidiano de vínculos e afetos. É um mundo que desfaz
lugares.

A questão da Palestina revela como a necessidade de um território
reconhecido pode tornar-se uma questão de sobrevivência coletiva, de
luta por soberania e de reconhecimento jurídico internacional. Gerações
de deslocamento forçado transformaram a demarcação territorial em
urgência existencial. Em Gaza e na Cisjordânia, o cotidiano é pautado
por um confinamento, onde o mundo se estreita entre postos de controle,
muros de concreto e a presença constante de um conflito que se tornou
paisagem. A ocupação militar e o bloqueio transformaram a própria ideia
de ``lugar'' --- como nos lembra Crespin, de espaço tecido por memória,
significado e identidade coletiva --- em algo sempre à beira do
desaparecimento. Quando um simples deslocamento exige autorização
militar, quando casas são demolidas sem aviso, quando comida, água,
eletricidade e cuidados médicos são controlados por uma força ocupante,
o mapa se desfaz. Resta a luta por um fragmento de chão que carregue não
só corpos, mas histórias. É nesse embate cotidiano, onde o viver é
resistência, que o território se revela como linguagem de pertencimento,
e o ato de permanecer, um gesto radical de afirmação cultural.
Infelizmente, a resistência cultural, por si só, não altera as
correlações de força e poder.

A Favela do Moinho enfrenta uma forma dura e brutal de deslocamento, por
meio da remoção forçada de famílias e da demolição de suas casas pelo
governo estadual, para promover o "desenvolvimento urbano". O uso de
força policial com gás lacrimogêneo e balas de borracha, e o
desmantelamento sistemático da comunidade, são uma tentativa de apagar o
seu "lugar" do mapa da cidade. Estão sistematicamente comprometendo a
integridade estrutural das casas restantes, criando riscos sanitários
com os escombros e cortando o acesso a serviços essenciais como
ambulâncias e transporte escolar.

Em escalas distintas, o efeito é semelhante nesses dois casos: a
fragilização do direito de permanecer.

É aqui que o livro de Crespin adquire uma urgência que vai além da
história judaica. Desde a Caxemira ao Saara Ocidental, das Ilhas Curilas
à região do Essequibo, até nossas próprias cidades, são pessoas que
construíram suas vidas e seu sentido de pertencimento num espaço que
outros grupos mais poderosos insistem em considerar ilegítimo. É preciso
interrogar o sistema que produz o exílio.

A questão não é apenas o direito a um lar, mas o direito de permanecer e
o direito de santificar e sacralizar este local. O direito de ter
raízes. De ser visto como alguém que pertence.

A troca comunitária emerge assim como única possibilidade de lugar. Não
um lugar onde a comunidade se reúne, mas o lugar que a comunidade cria
no próprio ato de reconhecer-se mutuamente. É um território móvel, feito
de gestos, palavras, cumplicidades --- um lugar que existe apenas
enquanto é praticado. E, assim, nos ensinam que há territorialidades que
não cabem no mapa. Que a identidade diaspórica e de nômades não é uma
condição de falta, mas uma forma específica e legítima de construir
identidade, um território construído através das migrações e comunidades
que se formam no trânsito. Como nos lembra Tim Ingold em "Contra o
espaço: lugar, movimento, conhecimento", o lugar não é um recipiente
para a vida, mas emerge do próprio movimento da vida. É no movimento, no
gesto de cozinhar da minha mãe e no ritual de colocar a mezuzá, que o
lugar se materializa.

Crespin nos lembra que a criação de um "lugar" é um ato de resistência
espiritual e cultural. Um exercício de sobrevivência que nos ensina a
navegar o sutil desequilíbrio entre aquilo que nos foi negado e a
coragem de reconhecer e reclamar o que legitimamente nos pertence. No
seguir habitando o mundo mesmo quando ele insiste em nos desabitar.

Mas resistir também implica nomear as estruturas que desmoronam o
``lugar''. Estados, alianças militares, projetos ideológicos e disputas
hegemônicas que instrumentalizam o território como ativo estratégico.
Reconhecer que mapas também são traçados por interesses geoeconômicos.

O livro conclui com a bênção \emph{Shehecheyanu
vekiyimanu vehigianu lazman haze}.

Hoje, esta bênção soa simultaneamente como oração, apelo e promessa.
Tranquiliza um pouco a dor, a raiva e a impotência que sinto neste
momento.

Uma oração por aqueles cujas vidas estão suspensas na geografia brutal
do conflito.

Um apelo para que possamos aprender a ver o mundo não como um espaço a
ser conquistado, mas como um lugar a ser partilhado.

Uma promessa, inspirada na resiliência diaspórica, de continuar a contar
as histórias, a manter vivas as fagulhas da memória e a insistir que
toda vida humana, em cada lugar contestado, merece encontrar seu próprio
mapa de volta para casa.

% \textbf{Comentário:}

% Acredito que seja excelente e condizente com a proposta do posfácio e
% também do trabalho de trazer as relações entre identidade,
% territorialidade, memória e conflito urbano/territorial. Importante
% trazer exemplos concretos sobre a condição de ter marcações territoriais
% tão definidas, como é o caso de Israel e Palestina --- justamente porque
% tento trazer no texto a ideia de que o judaísmo é uma cultura/identidade
% em diáspora, ou seja, nessa ideia de ausência de espaços tão marcados, o
% que constituiria, então o lugar?

% E o lugar é justamente esse movimento para dentro --- como a mezuzá, que é
% colocada nos batentes das portas no momento em que se entra em um
% cômodo, mais aprofundado dentro da casa --- então se coloca a mezuzá do
% lado direito durante a entrada do hall de apartamentos para o
% apartamento em si, da sala para o quarto\ldots{} e assim por diante.

% A experiência das comunidades em diáspora é sempre uma experiência da
% não espacialidade, a experiência do grupo, como uma experiência de
% peregrino, que não habita o espaço mas que cosmogoniza o espaço, o
% tranforma em lugar de apropriação e de pertencimento, até que esse mesmo
% lugar se desmanche e tenha que ser refeito em outro território.

% Essa ideia de território, ela pode ser tensionada tanto quanto um
% território real, mas também um território ficcional, um território
% identitário e comunitário. A construção de identidades diaspóricas que
% habitam o entre, entre-línguas, entre-códigos, entre-lugares.

% Acho bonita a ideia de comunidade, troca comunitária como lugar.

% Dando um feedback sobre o posfácio, acho que podemos trazer de outras
% formas a experiência identitária do \emph{estar entre}, como construção
% de identidade em si; usando de conceitos como da Gloria Anzanduia ou do
% capítulo `Contra o espaço: lugar, movimento, conhecimento' dentro do
% livro Estar Vivo do antropólogo inglês Tim Ingold.

\pagebreak